% !PRINTED 720
\weberpage{720}

%<<CONT
wodurch $D$ als eine von Null verschiedene positive Zahl definirt ist.
Was die Summen
\[
X = \overset{n}{\Sigma}\, \frac{\chi(n)}{n}
\]
betrifft, so ist daran zu erinnern, dass $\chi(n)$ [nach \S.\,193, (8)]
definirt ist durch
\weq{7}{n \equiv (-1)^\alpha 5^\beta \pmod{m}, \qquad
\chi(n) = (\pm 1)^\alpha \Theta^\beta,}
worin $\Theta$ eine $\nu^{\mathrm{te}}$ Einheitswurzel bedeutet.

Unter den $\nu^{\mathrm{ten}}$ Einheitswurzeln sind aber auch die
$\nu'^{\mathrm{ten}}$ Einheitswurzeln enthalten, und unter den
Charakteren $\chi$ auch die Charaktere f\"ur den K\"orper $\Omega_m'$.
Wir unterscheiden demnach \gesperrt{primitive und imprimitive Charaktere des
K\"orpers $\Omega_m$, indem wir die dem K\"orper $\Omega_m$
eigenth\"umlichen Charaktere $\chi$, d.\,h. die, in denen $\Theta$ eine
primitive $\nu^{\mathrm{te}}$ Einheitswurzel} ist, als primitiv
bezeichnen.

Unter den Summen $X$ kommen auch alle zu dem K\"orper $\Omega_m'$
geh\"origen Summen vor:
\[
X' = \overset{n}{\Sigma}\, \frac{\chi'(n)}{n}.
\]

Wenn man nun in (4) f\"ur $k, k'$ und $h_1, h_1'$ die im \S.\,195, (16),
(17) gebildeten Ausdr\"ucke einsetzt, und dann die Summen $X'$
weghebt, so ergiebt sich:
\weq{8}{\begin{aligned}
\pi^{\nu'} A &= 2 \cdot 2^{2^{\varkappa-4}(\varkappa-2)} \prod X_2 \\
D B &= 2^{2^{\varkappa-4}(\varkappa-2)} \prod X_1,
\end{aligned}}
\gesperrt{worin sich jetzt aber die Producte $\prod$ nur noch auf die mit den
primitiven Charakteren $\chi_1, \chi_2$ gebildeten Summen $X_1, X_2$
erstrecken.}

In den im \S.\,194, (12) gegebenen Ausdr\"ucken f\"ur die $X$:
\weq{9}{\begin{aligned}
X_1 &= -\frac{1}{2m} \sum_{1,\,m-1}^{s} f_1(r^s)
\log\left(\sin \frac{s\pi}{m}\right)^2 \\
X_2 &= -\frac{i\pi}{m^2} \sum_{1,\,m-1}^{s} s\, f_2(r^s)
\end{aligned}}
treten nun, unter der Voraussetzung, dass $m$ eine Potenz von 2 und
dass $\chi_1, \chi_2$ primitive Charaktere sind, bedeutende
Vereinfachungen ein.
